\section{남은 증명 의무와 권장 진행 순서}
\label{sec:open-obligations}

\subsection{Week~5/6이 닫은 실제 API 경계}

Week~5는 실제 caller의 control과 입력을, Week~6은 cross-call memory 조건을
다음 범위에서 닫았다.

\begin{itemize}
  \item \identifier{RawApiAddressBridge}가 canonical 정수 주소와
        \identifier{W64} 주소의 왕복을 증명한다.
  \item \identifier{keypair_raw_api_exports_matching_prefixes}가 실제 raw KeyGen의
        두 순차 export 뒤 외부 VK/SK의 첫 992바이트 일치를 증명한다.
  \item exact Sign/Verify trace, \(accept\Rightarrow tail\), accepted hash-input
        binding이 실제 caller의 control과 hash 인자를 노출한다.
  \item 실제 Sign의 signature/siglen 출력은 명시적으로 분리된 VK/SK/pre/message
        구간을 보존하고 KeyGen의 외부 key prefix를 유지한다.
  \item 실제 생성 hash는 전체 memory 등식 대신 VK/pre/message read region의
        바이트 등식만으로 비교된다.
  \item 실제 raw Sign 다음 cryptolab Verify의 실행과 두 trace의 순차 실행은
        결과와 최종 memory까지 정확히 대응한다.
\end{itemize}

\identifier{raw_api_region_reaches_mu_zero_loss}는 이 실제 external-memory 사실과
Sign import를 \identifier{sk_memory_prefix}로 바꾼다. 이 경로는
\cref{thm:memory-witness}의 constructed \identifier{stores}를 사용하지 않는다.

\subsection{동결된 보안 관문(gate): 저장 \texorpdfstring{\(\mu\)}{mu}의 생성/재전달(product/replay)}

이전의 \claim{OBL-API-KEY-MEMORY-RAW}는 구현의 early-reject 제어흐름보다 강한
무조건 명제였기 때문에 Week~5에서 분해되었다. Week~6 이후 직접 잔여 의무는
다음 하나의 의미론 간선으로 좁혀졌다.

\begin{obligation}[\claim{OBL-DIRECT-OBSERVED-MU}]
실제 raw Sign trace와 성공한 cryptolab Verify trace의 순차 실행 뒤
\[
  \operatorname{MuPrefix}_{32}
  (\mu_{S,\mathrm{trace}},\mu_{V,\mathrm{trace}})
\]
를 직접 결론으로 얻는다. 여기서 두 항은 각각
\identifier{SignRawApiMuTrace.observed_mu}와
\identifier{VerifyCryptolabMuTrace.observed_mu}다. 생성 hash return의 pRHL
정리를 이미 완료된 두 trace의 ghost field에 운반할 정당한 product/replay
규칙이 필요하다.
\end{obligation}

\cref{thm:sign-output-frame}과 \cref{thm:region-local-mu}, actual sequential trace는
frame, locality, control 및 input binding을 이미 공급한다. 하지만 Sign hash 뒤
signing continuation과 Verify hash 뒤 challenge/compare continuation이 서로
달라, 열린 signing losslessness 없이 한쪽 continuation을 버릴 수 없다. Week~7은
추가 observation wrapper, 목표를 반복하는 precondition, SHAKE locality axiom을
거부하고 \claim{OBL-MU-PRODUCT-REPLAY}를 \statusdeferred 인 명시적 paper
boundary로 동결했다. 따라서 이 lane에서
\(\delta_{\mu,\mathrm{raw\_api\_accept}}=0\)을 기록하지 않으며 같은 seam에
wrapper를 더 쌓지 않는다.

\subsection{최근 닫힌 관문: 실제 rANS 핵심과 production 전체 HBZ}

Week~7--11은 정준 서명 접두부, HBZ 잎/표, 순수 바이트 추적, 실제 접미부 복사와
결합 제어, 그리고 실제 인코더의 성공 접미부 정제를 차례로 닫았다. Week~12는
\claim{OBL-RANS-DECODE-REFINEMENT}를 정확 추적 부분정확성으로 닫았다. Week~13은
이 세 결과를 \identifier{Mode2RansActualHarness.run} 안에서 합성해
\claim{OBL-RANS-CORE-INVERSE}를 성공 조건부 Hoare 부분정확성으로 닫았다.
Week~14는 실제 full-HBZ prepare/apply와 이 핵심 역함수를 합성하고 production
SignaturePack/Unpack 경계까지 정확히 운반해
\claim{OBL-SIG-HBZ-ENCODE-DECODE}를 같은 강도의 부분정확성으로 닫았다.

\identifier{actual_rans_encode_copy_decode_inverse}의 전제는 초기 인자
바인딩과 \identifier{mode2_hbz_symbol_stream}뿐이다. 인코더가 실제로 실패하면
복사와 디코더를 실행하지 않고 초기 디코더 배열/상태를 유지한다. 인코더가
성공하면 실제 복사 구간에서 Week~12의 정확 추적 입력 전체를 구성하여
\[
\begin{aligned}
  bad_{\mathrm{dec}}&=0,&
  off_{\mathrm{dec}}&=\mathsf{encoded\_size},\\
  decoded[0..1024)&=expected[0..1024),&
  decoded[1024..2048)&=decoded_0[1024..2048)
\end{aligned}
\]
을 얻는다. 복사된 추적, 디코더 점별 읽기, 디코더 성공과 출력 등식은
사전조건이 아니다. \identifier{encoder_success_size_word_bridge}도 실제 성공
범위에서 W64 크기 뺄셈의 무랩을 도출하므로 임의의 정수 버퍼 증인을 두지
않는다.

\identifier{signature_pack_unpack_hbz_full_actual_exact}는 production 결합과
focused full-HBZ 결합의 반환 튜플 및 \identifier{decoder_ran}을 정확히
맞춘다. 그 Hoare 따름정리
\identifier{signature_pack_unpack_hbz_full_inverse_mode2}는 정준 HBZ만으로
zero-size 실패 또는 nonzero-size 성공 결론을 낸다. 성공 분기에서는
(4\le size\le1024), (bad=0), 1024계수 복원, coefficient/encoded 꼬리
프레임과 실제 prepare-symbol 추적을 얻는다. encoder success나 decoder 결과는
전제가 아니다.

\begin{verified}[\claim{OBL-RANS-ACTUAL-SUCCESS-WITNESS}]
Week~15의 \identifier{actual_rans_encode_all_six_success}는 실제 prepare가
만든 all-6 기호열에서 종료한 실제 encoder가 \(bad=0\), \(0\le off\le1020\),
\(4\le size\le1024\)와 정확한 trace/frame을 반환함을 증명한다.
\identifier{signature_pack_unpack_hbz_zero_success_mode2}는 이를 고정
all-zero HBZ의 production pack/unpack 왕복까지 올린다. 성공, 크기나 trace를
전제로 받지 않으므로 이 고정 입력 Hoare 의무는 \statusproved 다.
\end{verified}

Week~15 판정은 \identifier{GO-WITNESS}였지만
\identifier{GO-WITNESS-LOSSLESS}는 아니다. 고정 입력 실행의 종료,
losslessness, 확률 1 성공과 일반 정준 입력의 성공은 여전히 별도다. 이후
1주 최소 논문 범위가 원래 (h) 코덱 권고를 덮어썼으므로
\claim{OBL-SIG-H-ENCODE-DECODE}는 현재 \statusdeferred 다.

\subsection{KeyGen 종료 관문: STOP-KG-NTT}

Week~16은 실제 \identifier{_kp_m23_matrix}와
\identifier{_keypair_finalize_m23}를 투명한 두 호출 harness에서 합성했다.
\identifier{actual_m23_matrix_finalize_semantic_snapshot}는 실제 중간
\identifier{pre_bp}, NTT scratch, finalizer low/high 의미와 스냅샷
mod-\(2q\) 영 등식을 닫는다. 이 결과로 KG-2형 계수 분해와 KG-4의 조정된
\(s_2\) 성분은 증명되었지만 전체 논문 KeyGen 식은 아니다.

\begin{boundary}[\claim{STOP-KG-NTT}]
\identifier{KeygenM23ArithmeticSpec.output_row}가 표현하는
full-NTT--점별 곱--inverse-NTT 결과를 보안 모델의 행별 다항식 곱
\[
  \sum_{\mathrm{col}}
  A_{\mathrm{gen}}[\mathrm{row},\mathrm{col}]
  s_{\mathrm{gen}}[\mathrm{col}]
\]
과 연결하려면 odd-root orthogonality/full-NTT convolution 정리와
\identifier{Rq.poly}--security-list adapter가 필요하다. 둘 다 현재 checked
tree에 없으므로 faithful KG-1, KG-3, 전체 KG-4와
\(As=qj\pmod{2q}\)를 구성하지 않는다.
\end{boundary}

KeyGen 판정은 \identifier{STOP-KG-NTT}다. \claim{OBL-MINCORE-KEYGEN}은
KG-2/finalization에서 \statuspartial 로 동결되었다. MINCORE Sign은
\identifier{STOP-SIGN-CHAL-MODE2}, MINCORE Verify는
\identifier{STOP-VERIFY-MATRIX-CRT}이며 제한 합성은 \statusspecified 다.
Verify의 첫 재개 leaf는 \identifier{rq_mul_coeff_foldr_to_bigi}이며, 이어
2-by-4 NTT row-product와 CRT/freeze 절차 leaf가 남는다.

\subsection{분리된 기능정확성 경로}

\claim{OBL-SIGN-OUTPUT-TAIL-REACH}는 정당한 Sign 출력이 Verify의 unpack/norm
early-reject gate를 통과한다는 \statuspartial 의무다. signature prefix inverse는
닫혔지만 전체 suffix와 이후 gate는 열려 있다. 이는 성공한 Verify에
대해 tail 도달을 증명한 \cref{thm:verify-accept-tail}의 역이나 귀결이 아니다.
full pack/unpack inverse, norm, hint, response 및 산술 정확성이 필요하다. 이를
전체 challenge equality와 결합한 \claim{OBL-SIGN-VERIFY-CORRECTNESS}는 현재
\statusblocked 이다. 이 기능 경로를 accepted-forgery 보안 경로의 숨은 전제로
넣지 않는다.

\subsection{그 다음 관문: 문맥(context)과 전체 도전값 전사(challenge transcript)}

raw/internal 경로와 public-context 경로는 구현에서 실제로 다르다. Sign은
\identifier{_sf_prepare_pre_raw}와 \identifier{_sf_mu_preptr}를 사용하고,
Verify는 \identifier{prelen}의 최상위 비트로 context branch를 고른다. 다음을
보여야 한다.

\begin{enumerate}
  \item \(ctxlen\le255\)와 길이 인코딩의 정확성;
  \item 양쪽이 정확히 \(ctxlen\text{-byte}\Vert ctx\)를 같은 순서로 흡수함;
  \item caller memory가 두 절차의 pointer/length 전제를 공급함.
\end{enumerate}

그 다음 \claim{OBL-HIGH-LSB-PACK}은 Sign과 Verify가 만드는
\(\operatorname{HighBits}_h,\operatorname{LSB}\)의 packed bytes가 같음을
보여야 한다. 이 결과를 기존 \identifier{challenge_input_eq_components}와
합성하면 전체 challenge 입력 리스트가 같아진다. 이후에야 FIPS202
domain/padding/squeeze와 binary challenge sampler를 포함한 실제 full call을
paper ROM query에 연결할 수 있다.

\subsection{그 이후의 기능 및 확률적 의무}

\begin{longtable}{L{0.28\textwidth}L{0.18\textwidth}L{0.46\textwidth}}
\caption{현재 명명된 주요 미해결 의무}\label{tab:open-obligations}\\
\toprule
의무 & 성격 & 필요한 결과 \\
\midrule
\endfirsthead
\toprule
의무 & 성격 & 필요한 결과 \\
\midrule
\endhead
\claim{OBL-API-KEY-MEMORY-RAW-ACCEPT} & API/관계 & 승인된 Sign/Verify 추적
(accepted Sign/Verify trace)의 직접 저장 \(\mu_{32}\) 등식; 프레임/국소성/제어는
이미 증명됨 \\
\claim{OBL-DIRECT-OBSERVED-MU} & API/관계 & 실제 순차 추적 사후상태(actual
sequential-trace post-state)의 \identifier{observed_mu} 접두부 결론 \\
\claim{OBL-MU-PRODUCT-REPLAY} & 의미론/\statusdeferred & 생성 해시 반환값
(generated hash return)의 pRHL 정리를 완료된 추적 유령 필드(completed-trace
ghost field)로 옮기는 정당한 규칙 \\
\claim{OBL-SIG-PREFIX-FRAME} & 인코딩(encoding) & 패킹 서명 꼬리(pack signature
tail)와 도전값 출력 꼬리(challenge-output tail)의 보존; 언패킹 하위 꼬리(unpack
low tail)는 증명됨 \\
\claim{OBL-RANS-DECODE-REFINEMENT} & 인코딩(encoding);\newline\statusproved &
정확한 실제 추적, 상태 필드와 mode-2 표 아래 실제 디코더의 \(bad=0\), 정확한
소비 크기(consumed size), 1024기호 일치(symbol equality)와 출력 꼬리
프레임(frame); 종료는 별도 \\
\claim{OBL-RANS-ACTUAL-SUCCESS-WITNESS} & 인코딩(encoding);\newline\statusproved &
모든 기호가 6인 고정 입력에서 종료한 실제 인코더가 \(bad=0\), 정확한
크기/trace/frame을 반환함; 종료, losslessness와 \(\operatorname{phoare}=1\)은
별도 \\
\claim{OBL-RANS-CORE-INVERSE} & 인코딩(encoding);\newline\statusproved & 실제
인코더 실패 분기를 보존하고, 성공 시 실제 복사 구간을 디코더의 정확 추적
입력으로 운반해 \(bad=0\), 정확한 소비, 1024기호 복원과 꼬리 프레임을 얻는
성공 조건부 부분정확성; 종료/성공 도달은 별도 \\
\claim{OBL-SIG-HBZ-ENCODE-DECODE} & 인코딩(encoding);\newline\statusproved &
production 전체 HBZ 래퍼에서 zero-size 실패를 보존하고, nonzero-size 성공 시
크기 범위, \(bad=0\), 1024계수 복원, 두 꼬리 프레임과 prepare-symbol 추적을
공급하는 성공 조건부 부분정확성; 고정 zero 입력의 실패 배제는 증명됐고
종료/일반 입력 성공은 별도 \\
\claim{OBL-SIG-H-ENCODE-DECODE} & 인코딩(encoding);\newline\statusdeferred &
Week~16 최소 KeyGen 범위가 우선되어 연기된 실제 힌트 코덱 역함수 \\
\claim{OBL-MINCORE-KEYGEN} & 기능정확성;\newline
\statuspartial;\newline\identifier{STOP-KG-NTT} &
실제 snapshot/finalizer는 76대상 기준선에서 증명됨. 마지막
\identifier{output_row} NTT rewrite와 두 active row consequence는 77번째
파일에서 개별 컴파일됨. 그래도 odd-root convolution과
\identifier{Rq.poly}-to-security-list
adapter가 없어 faithful KG-1/KG-3/전체 KG-4와 논문 식은 동결 \\
\claim{OBL-MINCORE-SIGN} & 기능정확성;\newline\statuspartial;\newline
\identifier{STOP-SIGN-CHAL-MODE2} & 실제 round-challenge, z-check, hint helper
직접 호출과 accepted-branch control만 컴파일; full challenge semantics와
S-1--S-7 headline은 미증명 \\
\claim{OBL-MINCORE-VERIFY} & 기능정확성;\newline\statuspartial;\newline
\identifier{STOP-VERIFY-MATRIX-CRT} & actual helper 순서, V-1/V-2/V-5/V-6
machine-word 식, W64 norm gate와 tail/mismatch word expression은 컴파일;
matrix CRT 의미 bridge와 challenge equality는 미증명 \\
\claim{OBL-MINCORE-COMPOSITION} & 기능정확성;\newline\statusspecified &
decoded-object 경계의 제한 Sign--Verify 합성은 아직 미증명 \\
\claim{OBL-SIG-SUFFIX-METADATA},\newline
\claim{OBL-SIG-SUFFIX-PADDING} & 인코딩(encoding) & 크기 차이 왕복(size-delta
round trip)과 1474바이트 정준 영 패딩(canonical zero padding) \\
\claim{OBL-SIG-FULL-CODEC-MODE2} & 인코딩(encoding) & 접두부(prefix)와 전체
접미부(suffix)를 합친 실제 패킹/언패킹 역함수(actual pack/unpack inverse),
실패 표시(bad flag)의 성공과 정준 구문 분석(canonical parse) \\
\claim{OBL-SIGN-OUTPUT-TAIL-REACH} & 기능정확성 & 정당한 Sign 출력의 전체
패킹/언패킹(full pack/unpack) 및 노름/힌트 관문(norm/hint gate) 통과 \\
\claim{OBL-SIGN-VERIFY-CORRECTNESS} & 기능정확성 & 응답(response), 힌트(hint),
산술식과 전체 도전값을 포함한 정당한 서명의 승인(accept) \\
\claim{OBL-KG-TERMINATION}, \claim{OBL-KG-DIST} & 확률/종료 & 재시도 반복문
(retry loop)의 무손실성(losslessness) 또는 종료, 실제 KeyGen 출력분포, 특이성
거부(singular rejection) 효과 \\
\claim{OBL-FFT-TAIL}, TieMin & 수치해석 & 고정소수점 FFT 오차의 나쁜 사건
(bad event) 확률과 동률 최솟값(tied minimum) 정책의 차이 \\
\claim{OBL-SIGN-ACCEPTED-DIST} & 확률분포 & 초구 표본기(hyperball sampler),
거부(rejection), 승인 출력(accepted output)이 논문 분포와 일치하거나 명시적
통계손실을 가짐 \\
\claim{OBL-VERIFY-EVENT} & 대수/프로그램 & Jasmin의 최종 승인 비트(accept bit)가
논문의 노름(norm), 힌트, 방정식, 도전값 조건과 동치 \\
\claim{OBL-CHALLENGE} & 조합론/확률 & 도전값 지지집합(challenge support),
기수(cardinality), 점 확률(point probability), 최소 엔트로피(min-entropy) \\
\claim{OBL-FORK} & 암호 환원 & 두 새 전사(fresh transcript)에서 정확한
BST-MSIS 인스턴스(instance)를 추출 \\
\claim{OBL-SIGROM-ADAPTER} & 게임 기반 보안 & 바이트 수준 공용 API를 정확한
\identifier{Sig_ROM.Scheme} 및 질의 계수(query accounting)에 인스턴스화 \\
\bottomrule
\end{longtable}

\subsection{의존성에 따른 작업 순서}

현재 코드와 가장 잘 맞는 순서는 다음과 같다.

\begin{enumerate}
  \item 생성/재전달 경로는 현재 논문 경계로 유지하고 같은 호출자 이음매에
        관찰 래퍼를 더하지 않는다.
  \item Week~14 production 래퍼와 Week~15 고정 all-6 성공 정리를 유지하고
        rANS 반복 불변식을 다시 증명하지 않는다. 고정 입력 종료/losslessness는
        별도 의무로 남긴다.
  \item KG-NTT-MUL 감사의 \identifier{STOP-KG-NTT}를 유지한다. odd-root
        직교성/full-NTT convolution과 security-list adapter 없이는
        \identifier{pre_bp}를 \(A_{\mathrm{gen}}s_{\mathrm{gen}}\)으로 바꾸거나
        faithful KG-1/KG-3/전체 KG-4를 주장하지 않는다.
  \item \identifier{STOP-SIGN-CHAL-MODE2}를 유지한다. full
        highbits/LSB/\(\mu\)-to-SampleInBall leaf와 기존 convolution 의무 없이
        S-1--S-7을 주장하지 않는다.
  \item Verify의 첫 누락 leaf \identifier{rq_mul_coeff_foldr_to_bigi}를 닫고,
        2-by-4 NTT row-product와 CRT/freeze 절차 leaf를 거쳐 combined matrix
        theorem을 합성한 다음 downstream challenge leaf를 닫는다.
  \item (h) 코덱은 현재 연기 상태를 유지한다. 범위가 다시 열리면 크기
        메타데이터(size metadata), 영 패딩(zero padding), 패킹/도전값
        꼬리 프레임(pack/challenge tail frame)을 합성해 전체 1474바이트
        코덱(codec)과 정준 구문 분석(canonical parse)을 닫는다.
  \item 노름/힌트(norm/hint), 응답(response)과 산술 정확성을 결합하여
        \claim{OBL-SIGN-OUTPUT-TAIL-REACH}와 정당한 Sign/Verify 정확성
        (legitimate Sign/Verify correctness)을 닫는다.
  \item 별도의 보안 경로에서 공용 문맥(public-context) \(\mu\),
        highbits/LSB와 전체 도전값 호출(challenge call)을 합성한다.
  \item KeyGen/Sign의 확률분포와 종료 손실을 정량화하고 바이트 충실 ROM 스킴
        (byte-faithful ROM scheme)을 논문의 분포/포킹(distribution/forking)
        의무와 결합한다.
\end{enumerate}

이 순서는 넓은 보안 명제를 먼저 선언하고 하위 전제를 나중에 채우는 방식이
아니라, 이미 컴파일되는 가장 가까운 seam에서 바깥 경계로 한 단계씩 올라가는
방식이다.
